% ---
% title: Waiting for the next order
% tags: finance, exponential, density, distribution-function
% ---

\begin{problem}
	\begin{context}
		On an electronic exchange, buy and sell orders for a stock arrive at
		random moments. Traders and market makers care about how long they have
		to wait for the next order. The waiting time between orders is often
		modeled with an exponential distribution.
	\end{context}
	The waiting time $X$, in minutes, until the next order for a stock has
	density
	\[ f(x)=\begin{cases} 2e^{-2x}, & x\ge0,\\ 0, & x<0 \end{cases} \]
	(an exponential distribution with rate $\lambda=2$ per minute).
	(a) Check that $f$ is a density and find the distribution function $F$.
	(b) Find $P(X>1)$.
	(c) Find $P(X<0.5)$.
	(d) No order has arrived in the first half minute. What is the
	probability that we wait more than 1.5 minutes in total? Compare with (b).
\end{problem}

\begin{slideproblem}
	The waiting time $X$ (minutes) until the next order for a stock has
	density $f(x)=2e^{-2x}$ for $x\ge0$ (and $0$ for $x<0$).
	(a) Check that $f$ is a density; find $F$.
	\ (b) $P(X>1)$.
	\ (c) $P(X<0.5)$.
	\ (d) No order in the first half minute. Probability that we wait more
	than 1.5 minutes in total? Compare with (b).
\end{slideproblem}

\begin{slides}
	\small
	\exstatement
	\pause
	\begin{columns}[T]
		\begin{column}{.58\textwidth}
			\sol (a) $f(x)\ge0$, and
			\[ \int_0^{\infty}2e^{-2x}\,dx=\Bigl[-e^{-2x}\Bigr]_0^{\infty}
				=0-(-1)=1. \]
			\pause
			For $x\ge0$:
			\[ F(x)=\int_0^{x}2e^{-2t}\,dt=\Bigl[-e^{-2t}\Bigr]_0^{x}
				=1-e^{-2x}, \]
			$F(x)=0$ for $x<0$.
			\pause
			\ Tail: \ $\alert{P(X>x)=e^{-2x}}$.
		\end{column}
		\begin{column}{.40\textwidth}
			\centering
			\begin{tikzpicture}[xscale=1.5,yscale=1.1]
				\fill[teal!30] (1,0) -- plot[domain=1:2.4,samples=40]
				(\x,{2*exp(-2*\x)}) -- (2.4,0) -- cycle;
				\draw[->] (-.2,0) -- (2.6,0) node[right] {\scriptsize $x$};
				\draw[->] (0,-.1) -- (0,2.3) node[above] {\scriptsize $f(x)$};
				\draw[navy,thick] plot[domain=0:2.4,samples=50]
				(\x,{2*exp(-2*\x)});
				\foreach \x in {0.5,1,1.5,2}
				\draw (\x,.04) -- (\x,-.04) node[below] {\scriptsize $\x$};
				\draw (.04,2) -- (-.04,2) node[left] {\scriptsize $2$};
				\draw[teal,->] (1.6,.7) node[above] {\scriptsize $P(X>1)$}
				-- (1.2,.12);
			\end{tikzpicture}
		\end{column}
	\end{columns}
\end{slides}

\begin{slides}
	\small
	(b) $P(X>1)=e^{-2\cdot1}=e^{-2}\approx0.135$.\\[2mm]
	\pause
	(c) $X$ is continuous, so $P(X=0.5)=0$ and
	\[ P(X<0.5)=F(0.5)=1-e^{-2\cdot0.5}=1-e^{-1}\approx0.632. \]
	\pause
	(d) ``More than 1.5 in total'' already contains ``more than 0.5'':
	\[ P(X>1.5\mid X>0.5)=\frac{P(X>1.5)}{P(X>0.5)}
		=\frac{e^{-3}}{e^{-1}}=e^{-2}\approx0.135. \]
	\pause
	The same as (b)! Having waited half a minute does not change the chance of
	waiting more than one more minute.\\[1mm]
	\alert{The exponential distribution has no memory.}
\end{slides}

\begin{solution}
	\idea A density must be nonnegative and integrate to 1. The distribution
	function is the integral of the density from the left up to $x$. Then
	every probability about $X$ comes from $F$; for (d) we use the definition
	of conditional probability.

	\step{(a)} For $x\ge0$, $2e^{-2x}>0$ because an exponential is always positive,
	and $f(x)=0$ for $x<0$, so $f(x)\ge0$ everywhere. The total area under $f$ is
	\[ \int_{-\infty}^{\infty}f(x)\,dx=\int_0^{\infty}2e^{-2x}\,dx. \]
	An antiderivative of $2e^{-2x}$ is $-e^{-2x}$ (differentiate: the chain
	rule gives $-e^{-2x}\cdot(-2)=2e^{-2x}$). At the upper bound,
	$-e^{-2x}\to0$ as $x\to\infty$; at the lower bound it equals $-e^{0}=-1$.
	So
	\[ \int_0^{\infty}2e^{-2x}\,dx=\Bigl[-e^{-2x}\Bigr]_0^{\infty}
		=0-(-1)=1, \]
	and $f$ is a density.

	For the distribution function, $F(x)=\int_{-\infty}^{x}f(t)\,dt$. If
	$x<0$, the density is 0 to the left of $x$, so $F(x)=0$. If $x\ge0$, only
	the part from 0 to $x$ contributes:
	\[ F(x)=\int_0^{x}2e^{-2t}\,dt=\Bigl[-e^{-2t}\Bigr]_0^{x}
		=-e^{-2x}-(-1)=1-e^{-2x}. \]
	So
	\[ F(x)=\begin{cases} 0, & x<0,\\ 1-e^{-2x}, & x\ge0. \end{cases} \]
	A useful consequence is the \emph{tail probability}: for $x\ge0$,
	\[ P(X>x)=1-F(x)=e^{-2x}. \]
	\begin{center}
		\begin{tikzpicture}[xscale=2,yscale=1.2]
			\fill[teal!30] (1,0) -- plot[domain=1:2.6,samples=40]
			(\x,{2*exp(-2*\x)}) -- (2.6,0) -- cycle;
			\draw[->] (-.2,0) -- (2.8,0) node[right] {\small $x$};
			\draw[->] (0,-.1) -- (0,2.3) node[above] {\small $f(x)$};
			\draw[navy,thick] plot[domain=0:2.6,samples=60] (\x,{2*exp(-2*\x)});
			\foreach \x in {0.5,1,1.5,2,2.5}
			\draw (\x,.04) -- (\x,-.04) node[below] {\scriptsize $\x$};
			\draw (.03,2) -- (-.03,2) node[left] {\scriptsize $2$};
			\draw[teal,->] (1.6,.8) node[above] {\scriptsize area $=P(X>1)=e^{-2}$}
			-- (1.2,.12);
		\end{tikzpicture}
	\end{center}

	\step{(b)} Use the tail formula with $x=1$:
	\[ P(X>1)=e^{-2\cdot1}=e^{-2}\approx0.135. \]
	This is the shaded area in the picture. There is about a 13.5\% chance
	that no order arrives within a minute.

	\step{(c)} For a continuous random variable a single point has
	probability 0, so $P(X<0.5)=P(X\le0.5)=F(0.5)$. Hence
	\[ P(X<0.5)=1-e^{-2\cdot0.5}=1-e^{-1}\approx1-0.368=0.632. \]

	\step{(d)} We want $P(X>1.5\mid X>0.5)$. By definition,
	\[ P(X>1.5\mid X>0.5)=\frac{P(\{X>1.5\}\cap\{X>0.5\})}{P(X>0.5)}. \]
	If $X>1.5$, then certainly $X>0.5$, so the intersection is just the event
	$\{X>1.5\}$. Using the tail formula twice,
	\[ P(X>1.5\mid X>0.5)=\frac{P(X>1.5)}{P(X>0.5)}
		=\frac{e^{-2\cdot1.5}}{e^{-2\cdot0.5}}=\frac{e^{-3}}{e^{-1}}
		=e^{-2}\approx0.135. \]
	This is exactly the answer of (b). Having already waited half a minute,
	the chance of waiting more than one more minute is the same as the chance
	of waiting more than one minute from the start. The past waiting does not
	matter: the exponential distribution is \emph{memoryless}. In general,
	$P(X>s+t\mid X>s)=e^{-\lambda(s+t)}/e^{-\lambda s}=e^{-\lambda t}=P(X>t)$.
	\begin{answer}
		(a) $\int_0^\infty 2e^{-2x}\,dx=1$; $F(x)=1-e^{-2x}$ for $x\ge0$,
		$F(x)=0$ for $x<0$
		\quad (b) $e^{-2}\approx0.135$
		\quad (c) $1-e^{-1}\approx0.632$
		\quad (d) $e^{-2}\approx0.135$, the same as (b)
	\end{answer}
	\begin{remark}
		\textbf{Common mistake:} in (d), computing
		$P(X>1.5)=e^{-3}\approx0.050$ and ignoring the information that no
		order came in the first half minute; or computing
		$P(0.5<X\le1.5)$. The condition shrinks the sample space to
		$\{X>0.5\}$, so we divide by its probability. In finance terms: if
		orders arrive in a memoryless way, a quiet market in the last half
		minute tells us nothing about how much longer the quiet will last.
	\end{remark}
\end{solution}
